% Baccalauréat malgache, série D, 2007 : corrigé.
\TCDBacAttribution{D}{2007}
\providecommand{\Plan}{(\mathscr{P})}
\providecommand{\Ouv}{(O ; \vec{u} , \vec{v})}
\providecommand{\Courbe}{(\mathscr{C})}
\providecommand{\cm}{~\mathrm{cm}}
\providecommand{\dx}{~\mathrm{d}x}
%
%%%%%%%%%%%%%%%%%%%%%%%%%%%%%%%%%%%%%%%%%%%%%%%%%%%%%%%%%%%%%%%%%%%%%%%%%%%%%%%%%%%%%%%%%%%%%%%%%%%%%%%%%%%%%%%%%%%%%%%%%%%%%%%%
\noindent \underline{\textsc{\textbf{Exercice I}}} (5 points)
\begin{enumerate}
   \item
      \begin{enumerate}
         \item \vln{Calcul de $P(-3i)$ :}
            $(-3i)^3 = 27i$, $(1 + i)(-3i)^2 = -9 - 9i$, $(4 - i)(-3i) = -3 - 12i$. Donc
            $P(-3i) = 27i - 9 - 9i - 3 - 12i - 6i + 12$, \ansbox{$P(-3i) = 0$.}
         \item \vln{Racine r\'eelle :}
            Pour $x$ r\'eel, $P(x) = (x^3 + x^2 + 4x + 12) + i\,(x^2 - x - 6)$. $P(x) = 0$ exige
            $x^2 - x - 6 = 0$, soit $x = 3$ ou $x = -2$, et $x^3 + x^2 + 4x + 12 = 0$.
            Pour $x = 3$ : $27 + 9 + 12 + 12 = 60 \neq 0$ ; pour $x = -2$ : $-8 + 4 - 8 + 12 = 0$.
            \par\noindent\ansbox{La racine r\'eelle est $\alpha = -2$.}
         \item \vln{Calcul de $\beta$ :}
            En d\'eveloppant, $(z + 3i)(z + 2)(z + \beta)$ a pour terme constant $6i\beta$ et pour coefficient de $z^2$ : $2 + 3i + \beta$.
            Par identification : $2 + 3i + \beta = 1 + i$ donne $\beta = -1 - 2i$, et l'on v\'erifie $6i\beta = -6i + 12$
            ainsi que le coefficient de $z$ : $6i + (2 + 3i)\beta = 6i + (4 - 7i) = 4 - i$.
            \par\noindent\ansbox{$\beta = -1 - 2i$ : $P(z) = (z + 3i)(z + 2)(z - 1 - 2i)$.}
         \item \vln{Solutions :} \ansbox{$S_{\mathbb{C}} = \{ -3i \,;\, -2 \,;\, 1 + 2i \}$.}
      \end{enumerate}
   \item \vln{Angle, rapport et nature de $ABC$ :}
      (Dans l'\'enonc\'e, $Z_B$ d\'esigne $z_B$.)
      \[
         \frac{z_C - z_B}{z_A - z_B} = \frac{3 + 2i}{2 - 3i} = \frac{(3 + 2i)(2 + 3i)}{13} = \frac{13i}{13} = i.
      \]
      Donc $\left|\frac{z_C - z_B}{z_A - z_B}\right| = 1$ et $\left(\overrightarrow{BA} ; \overrightarrow{BC}\right) = \arg(i) = \frac{\pi}{2}$ $[2\pi]$.
      \par\noindent\ansbox{$BC = BA$ et $\left(\overrightarrow{BA} ; \overrightarrow{BC}\right) = \frac{\pi}{2}$ : $ABC$ est rectangle isoc\`ele en $B$.}
   \item \vln{Expression complexe de $r$ :}
      $z' - z_B = \expo{i\frac{\pi}{2}}(z - z_B)$, soit $z' = i(z + 2) - 2$ :
      \par\noindent\ansbox{$r : z' = iz - 2 + 2i$} (et $r(A) = C$ d'apr\`es la question 2).
   \item
      \begin{enumerate}
         \item \vln{Nature de $S$ :}
            $S : z' = 2iz - 2 + 4i$ avec $2i = 2\expo{i\frac{\pi}{2}}$ : rapport $2$, angle $\frac{\pi}{2}$.
            Centre : $\omega = \dfrac{-2 + 4i}{1 - 2i} = \dfrac{(-2 + 4i)(1 + 2i)}{5} = \dfrac{-10}{5} = -2$.
            \par\noindent\ansbox{$S$ est la similitude directe de centre $B(-2 ; 0)$, de rapport $2$ et d'angle $\frac{\pi}{2}$.}
         \item \vln{Transformation $h$ telle que $r \circ h = S$ :}
            $r$ est bijective, donc $h = r^{-1} \circ S$. $r^{-1}$ est la rotation de centre $B$ et d'angle $-\frac{\pi}{2}$ :
            $h$ est une similitude directe de centre $B$ (fix\'e par $S$ et par $r^{-1}$), de rapport $2 \times 1 = 2$ et
            d'angle $\frac{\pi}{2} - \frac{\pi}{2} = 0$ : c'est une homoth\'etie.
            En complexes : $r^{-1} : z' = -i(z + 2) - 2$, donc $h : z' = -i(2iz - 2 + 4i + 2) - 2 = 2z + 2$.
            \par\noindent\ansbox{$h$ est l'homoth\'etie de centre $B(-2 ; 0)$ et de rapport $2$ : $z' = 2z + 2$.}
            (Contr\^ole : $r(h(z)) = i(2z + 4) - 2 = 2iz - 2 + 4i$.)
      \end{enumerate}
\end{enumerate}
%%%%%%%%%%%%%%%%%%%%%%%%%%%%%%%%%%%%%%%%%%%%%%%%%%%%%%%%%%%%%%%%%%%%%%%%%%%%%%%%%%%%%%%%%%%%%%%%%%%%%%%%%%%%%%%%%%%%%%%%%%%%%%%%
\noindent \underline{\textsc{\textbf{Exercice II}}} (5 points) \newline
D'apr\`es le tableau, les six jetons sont
A (jaune), A (noir), B (rouge), B (noir), C (rouge), D (jaune).
Il y a donc 2 jetons rouges, 2 jetons portant une voyelle (les deux A) et 4 jetons portant une consonne.
\begin{description}
   \item[I.] Tirage successif sans remise de trois jetons : $6 \times 5 \times 4 = 120$ tirages ordonn\'es \'equiprobables.
      \begin{enumerate}
         \item \vln{Probabilit\'es de $E$ et $F$ :}
            $E$ : on choisit la place du jeton rouge (3 fa\c{c}ons), ce jeton (2 choix), puis deux jetons non rouges ordonn\'es
            ($4 \times 3 = 12$) : $3 \times 2 \times 12 = 72$ tirages. \ansbox{$P(E) = \frac{72}{120} = \frac{3}{5}$.}

            $F$ : un B (2 choix), puis un A (2 choix), puis le C (1 choix) : $4$ tirages. \ansbox{$P(F) = \frac{4}{120} = \frac{1}{30}$.}
         \item
            \begin{enumerate}
               \item \vln{Univers image de $X$ :}
                  Il n'y a que deux voyelles : au plus tard le troisi\`eme jeton est une consonne. \ansbox{$X(\Omega) = \{1 ; 2 ; 3\}$.}
               \item \vln{Loi et esp\'erance de $X$ :}
                  $P(X = 1) = \frac{4}{6} = \frac{2}{3}$ ; $P(X = 2) = \frac{2}{6}\times\frac{4}{5} = \frac{4}{15}$ ;
                  $P(X = 3) = \frac{2}{6}\times\frac{1}{5}\times\frac{4}{4} = \frac{1}{15}$.
                  \[
                     \begin{array}{|c|c|c|c|}
                        \hline
                        k & 1 & 2 & 3 \\ \hline
                        P(X = k) & \frac{2}{3} & \frac{4}{15} & \frac{1}{15} \\ \hline
                     \end{array}
                  \]
                  $E(X) = \frac{10}{15} + \frac{8}{15} + \frac{3}{15} = \frac{21}{15}$, \ansbox{$E(X) = \frac{7}{5}$.}
            \end{enumerate}
      \end{enumerate}
   \item[II.] Tirage simultan\'e : $\mathrm{C}_6^3 = 20$ tirages \'equiprobables.
      \begin{enumerate}
         \item \vln{Probabilit\'e de $G$ :}
            Il faut un B (2 choix), un A (2 choix) et le C : $4$ tirages. \ansbox{$P(G) = \frac{4}{20} = \frac{1}{5}$.}
         \item
            \begin{enumerate}
               \item \vln{Loi de $Y$ :}
                  Les trois \'epreuves sont ind\'ependantes et $G$ a la probabilit\'e $\frac{1}{5}$ \`a chacune :
                  $Y$ suit la loi binomiale $\mathscr{B}\left(3 ; \frac{1}{5}\right)$, $Y(\Omega) = \{0 ; 1 ; 2 ; 3\}$ et
                  $P(Y = k) = \mathrm{C}_3^k \left(\frac{1}{5}\right)^k\left(\frac{4}{5}\right)^{3-k}$.
                  \[
                     \begin{array}{|c|c|c|c|c|}
                        \hline
                        k & 0 & 1 & 2 & 3 \\ \hline
                        P(Y = k) & \frac{64}{125} & \frac{48}{125} & \frac{12}{125} & \frac{1}{125} \\ \hline
                     \end{array}
                  \]
               \item \vln{Esp\'erance et variance :}
                  \par\noindent\ansbox{$E(Y) = 3 \times \frac{1}{5} = \frac{3}{5}$ et $V(Y) = 3 \times \frac{1}{5} \times \frac{4}{5} = \frac{12}{25}$.}
            \end{enumerate}
      \end{enumerate}
\end{description}
%%%%%%%%%%%%%%%%%%%%%%%%%%%%%%%%%%%%%%%%%%%%%%%%%%%%%%%%%%%%%%%%%%%%%%%%%%%%%%%%%%%%%%%%%%%%%%%%%%%%%%%%%%%%%%%%%%%%%%%%%%%%%%%%
\noindent \underline{\textsc{\textbf{Probl\`eme}}} (10 points)
\begin{description}
   \item[I.] $g(x) = 1 - (x - 1)^2\expo{-x}$.
      \begin{description}
         \item[a)] \vln{Limite en $-\infty$ :}
            $(x - 1)^2 \to +\infty$ et $\expo{-x} \to +\infty$, \ansbox{$\lim\limits_{x\to-\infty} g(x) = -\infty$.}
         \item[b)] \vln{Limite en $+\infty$ :}
            $(x - 1)^2\expo{-x} = \dfrac{(x-1)^2}{x^2}\cdot\dfrac{x^2}{\expo{x}}$ ; le premier facteur tend vers $1$ et le second vers $0$
            (car $\frac{\expo{x}}{x^2} \to +\infty$). \ansbox{$\lim\limits_{x\to+\infty} g(x) = 1$} : la droite $y = 1$ est asymptote.
         \item[c)] \vln{D\'eriv\'ee :}
            $g'(x) = -\left[2(x - 1)\expo{-x} - (x - 1)^2\expo{-x}\right] = -(x - 1)\expo{-x}(2 - x + 1)$,
            \par\noindent\ansbox{$g'(x) = (x - 1)(x - 3)\,\expo{-x}$.}
         \item[d)] \vln{Tableau de variation :}
            $\expo{-x} > 0$ ; $g'(x)$ a le signe du trin\^ome $(x - 1)(x - 3)$. $g(1) = 1$, $g(3) = 1 - 4\expo{-3} \simeq 0{,}80$.
            \[
               \begin{array}{c|ccccccc}
                  x & -\infty & & 1 & & 3 & & +\infty \\ \hline
                  g'(x) & & + & 0 & - & 0 & + & \\ \hline
                  & & & 1 & & & & 1 \\
                  g & & \nearrow & & \searrow & & \nearrow & \\
                  & -\infty & & & & 1 - 4\expo{-3} & &
               \end{array}
            \]
         \item[e)] \vln{Signe de $g$ :}
            $g(0) = 1 - 1 = 0$. Sur $]-\infty ; 1]$, $g$ est strictement croissante et s'annule en $0$ : $g < 0$ sur $]-\infty ; 0[$,
            $g > 0$ sur $]0 ; 1]$. Sur $[1 ; +\infty[$, $g(x) \geq \min\left(g(3) ; 1\right) = 1 - 4\expo{-3} > 0$.
            \par\noindent\ansbox{$g(x) < 0$ si $x < 0$ ; $g(0) = 0$ ; $g(x) > 0$ si $x > 0$.}
      \end{description}
   \item[II.]
      \begin{enumerate}
         \item
            \begin{enumerate}
               \item \vln{Limites de $f$ :}
                  En $-\infty$ : $f(x) = \expo{-x}\left(x^2 + 1 + x\expo{x}\right)$, avec $x\expo{x} \to 0$ et $x^2 + 1 \to +\infty$,
                  $\expo{-x} \to +\infty$ : \ansbox{$\lim\limits_{x\to-\infty} f(x) = +\infty$.}
                  En $+\infty$ : $(x^2 + 1)\expo{-x} = \frac{x^2}{\expo{x}} + \expo{-x} \to 0$ et $x \to +\infty$ : \ansbox{$\lim\limits_{x\to+\infty} f(x) = +\infty$.}
               \item \vln{$f' = g$ et variations :}
                  $f'(x) = 2x\expo{-x} - (x^2 + 1)\expo{-x} + 1 = 1 - (x^2 - 2x + 1)\expo{-x} = 1 - (x - 1)^2\expo{-x}$,
                  \par\noindent\ansbox{$f'(x) = g(x)$.}
                  \[
                     \begin{array}{c|ccccc}
                        x & -\infty & & 0 & & +\infty \\ \hline
                        f'(x) & & - & 0 & + & \\ \hline
                        & +\infty & & & & +\infty \\
                        f & & \searrow & & \nearrow & \\
                        & & & 1 & &
                     \end{array}
                  \]
                  Minimum : $f(0) = 1$.
            \end{enumerate}
         \item
            \begin{enumerate}
               \item \vln{Point d'inflexion en $x = 1$ :}
                  $f''(x) = g'(x) = (x - 1)(x - 3)\expo{-x}$ s'annule en $1$ en changeant de signe (positive avant $1$,
                  n\'egative entre $1$ et $3$). \ansbox{$A\left(1 ; 1 + \frac{2}{e}\right)$ est un point d'inflexion de $\Courbe$}
                  ($f(1) = 2\expo{-1} + 1 \simeq 1{,}74$).
               \item \vln{Tangente en $A$ :}
                  $f'(1) = g(1) = 1$ : $y = (x - 1) + 1 + \frac{2}{e}$, \ansbox{$(T) : y = x + \frac{2}{e}$.}
               \item \vln{Asymptote $(\Delta)$ :}
                  $f(x) - x = (x^2 + 1)\expo{-x} \to 0$ en $+\infty$ : \ansbox{$(\Delta) : y = x$ est asymptote oblique \`a $\Courbe$ en $+\infty$.}
                  Comme $(x^2 + 1)\expo{-x} > 0$, $\Courbe$ est toujours au-dessus de $(\Delta)$.
               \item \vln{Branche infinie en $-\infty$ :}
                  $\dfrac{f(x)}{x} = \dfrac{x^2 + 1}{x}\,\expo{-x} + 1$ ; $\dfrac{x^2 + 1}{x} \to -\infty$ et $\expo{-x} \to +\infty$, donc
                  $\dfrac{f(x)}{x} \to -\infty$.\par\noindent\ansbox{En $-\infty$, $\Courbe$ admet une branche parabolique de direction $(Oy)$.}
               \item \vln{Trac\'e :} (quelques valeurs : $f(-1) \simeq 4{,}44$, $f(2) \simeq 2{,}68$, $f(3) \simeq 3{,}50$)
                  \begin{center}
                  \begin{tikzpicture}[scale=0.7,>=stealth]
                     \draw[very thin,gray!30] (-2,-1) grid (8,9);
                     \draw[->] (-2,0) -- (8,0) node[below left] {$x$};
                     \draw[->] (0,-1) -- (0,9) node[below left] {$y$};
                     \foreach \t in {-1,1,2,3,4,5,6,7} \draw (\t,0.08) -- (\t,-0.08) node[below] {\scriptsize $\t$};
                     \foreach \t in {1,2,...,8} \draw (0.08,\t) -- (-0.08,\t) node[left] {\scriptsize $\t$};
                     \begin{scope}
                        \clip (-2,-1) rectangle (8,9);
                        \draw[loko4,thick] plot[domain=-1:8] (\x,\x);
                        \draw[loko3,thick,dashed] plot[domain=-2:8] (\x,{\x+0.7358});
                        \draw[loko1,very thick,smooth,samples=150] plot[domain=-1.35:8] (\x,{(\x*\x+1)*exp(-\x)+\x});
                     \end{scope}
                     \fill (1,1.7358) circle (2pt) node[below right] {$A$};
                     \node[loko4] at (7.5,6.8) {$(\Delta)$};
                     \node[loko3] at (6.6,8.2) {$(T)$};
                     \node[loko1] at (-0.8,7.5) {$\Courbe$};
                  \end{tikzpicture}
                  \end{center}
            \end{enumerate}
         \item
            \begin{enumerate}
               \item \vln{Calcul de $I$ :}
                  Premi\`ere int\'egration par parties ($u = x^2$, $v' = \expo{-x}$, $u' = 2x$, $v = -\expo{-x}$) :
                  $I = \left[-x^2\expo{-x}\right]_0^1 + 2\int_0^1 x\expo{-x}\dx = -\expo{-1} + 2J$.
                  Seconde ($u = x$, $v' = \expo{-x}$) : $J = \left[-x\expo{-x}\right]_0^1 + \int_0^1 \expo{-x}\dx = -\expo{-1} + 1 - \expo{-1} = 1 - 2\expo{-1}$.
                  Donc $I = -\expo{-1} + 2 - 4\expo{-1}$, \ansbox{$I = 2 - \frac{5}{e}$} ($\simeq 0{,}16$).
               \item \vln{Aire $A$ :}
                  Sur $[0 ; 1]$, $\Courbe$ est au-dessus de $(\Delta)$ ; l'\'ecart vaut $(x^2 + 1)\expo{-x}$. L'unit\'e d'aire vaut $1\cm^2$.
                  \[
                     A = \int_0^1 (x^2 + 1)\expo{-x}\dx = I + \Big[-\expo{-x}\Big]_0^1 = 2 - \frac{5}{e} + 1 - \frac{1}{e},
                  \]
                  \par\noindent\ansbox{$A = \left(3 - \frac{6}{e}\right)\cm^2 \simeq 0{,}79\cm^2$.}
                  (L'arrondi $\expo{-1} \simeq 0{,}4$ propos\'e par l'\'enonc\'e est grossier et donnerait $0{,}6\cm^2$ ; avec $\expo{-1} \simeq 0{,}368$ on obtient $0{,}79\cm^2$.)
            \end{enumerate}
      \end{enumerate}
\end{description}
